\section{System Overview}
\label{sec:system}

\subsection{Design Rationale: Decoupled Algorithm and Simulation Layers}

The implementation is organised around a strict two-layer separation.
The \emph{algorithm layer} contains all learning components: a Gaussian-process
dynamics model, Monte Carlo policy gradient optimiser, radial-basis-function (RBF)
speed policy, and a ballistic simulator implemented in NumPy and PyTorch.
The \emph{simulation/visualisation layer} is either (a) the same NumPy ballistic
simulator used for training or (b) a PyBullet scene with a KUKA iiwa7 arm for
validation and demonstration.
These two layers share a single interface: the commanded release velocity
$\mathbf{v}_{\mathrm{cmd}} \in \mathbb{R}^3$.

The central design decision is that the arm's motion is \emph{cosmetic}.
Ball velocity at release is set programmatically via PyBullet's
\texttt{resetBaseVelocity} to exactly $\mathbf{v}_{\mathrm{cmd}}$, or
$\mathbf{v}_{\mathrm{cmd}} + \boldsymbol{\epsilon}$ when sensor noise is enabled.
Arm joint trajectories are planned and executed for visual fidelity, but any
joint-tracking error has no effect on the ball's actual initial velocity.
Consequently, GP training data are never corrupted by kinematic imprecision, and the
algorithm's data-efficiency guarantees from \cite{turcato2025mcpilot} transfer intact
to the visualisation environment.

\subsection{Sub-Project Structure}

The repository is organised into four sub-projects, each adding one independent
experimental variable over the previous:

\begin{itemize}
  \item \textbf{\texttt{mc-pilot/}} — Baseline NumPy reproduction of \mcpilot{}
        \cite{turcato2025mcpilot}.  Five targets on a flat plane ($z_{\mathrm{release}}
        = 0$\,m), no sensor noise.  All four failure modes were diagnosed and fixed
        here (Section~\ref{sec:implementation}).

  \item \textbf{\texttt{mc-pilot-elevated/}} — Release-height study.  Sweeps five
        release heights $z \in \{0.0, 0.5, 1.0, 1.5, 2.0\}$\,m using the NumPy
        ballistic simulator; arm geometry is not modelled.  Isolates the effect of
        geometric variation on GP generalisation and landing-error distribution.

  \item \textbf{\texttt{mc-pilot-pybullet/}} — PyBullet arm integration with realistic
        noise models.  Introduces velocity-slip noise
        ($\boldsymbol{\epsilon} \sim \mathcal{N}(\mathbf{0}, \sigma^2 I)$) injected at
        release.  Studies the zero-mean-noise unlearnability regime
        (Section~\ref{sec:noise}).

  \item \textbf{\texttt{mc-pilot-pb-elevated/}} — Combined PyBullet scene with elevated
        release positions.  Used for end-to-end visualisation and qualitative validation
        only; no new algorithm experiments are run here.
\end{itemize}

\subsection{Ballistic Simulator}

The NumPy ballistic simulator integrates the 6D ball state
$\mathbf{s} = [x, y, z, v_x, v_y, v_z]^\top$ under gravity and aerodynamic drag.
Drag follows Equation~35 of \cite{turcato2025mcpilot} with a
Reynolds-number-dependent drag coefficient $C_D$.
The state update uses the paper's reconstruction rule (Eq.~18):
\begin{align}
  \mathbf{v}_{t+1} &= \mathbf{v}_t + \Delta\mathbf{v}_t, \label{eq:vel_update} \\
  \mathbf{p}_{t+1} &= \mathbf{p}_t + T_s\,\mathbf{v}_t + \tfrac{T_s}{2}\,\Delta\mathbf{v}_t,
  \label{eq:pos_update}
\end{align}
where $T_s = 0.02$\,s is the integration timestep.
The 6D state is augmented to 8D by appending the target coordinates
$[P_x, P_y]$, forming the input to the GP dynamics model.
Landing is detected at $z \leq 0$ and the exact crossing is resolved by linear
interpolation between the last two timesteps, avoiding systematic over-shoot error.

\subsection{PyBullet Arm Pipeline}

The \texttt{ArmController} class wraps a KUKA iiwa7 URDF (loaded from
\texttt{pybullet\_data}) mounted at a configurable base position.  The throw
execution proceeds in five steps:

\begin{enumerate}
  \item \textbf{IK for release pose.}  The desired release Cartesian position is
        converted to joint angles via PyBullet's built-in IK solver.

  \item \textbf{Joint velocity mapping.}  A Jacobian pseudoinverse maps the desired
        end-effector velocity to joint velocities, subject to scaled iiwa7 hardware
        limits (\texttt{vel\_limit\_multiplier} parameter).

  \item \textbf{Three-phase cubic trajectory.}  A piecewise cubic spline drives the arm
        through neutral $\to$ windup $\to$ release $\to$ follow-through waypoints.

  \item \textbf{Ball attachment and release.}  A \texttt{JOINT\_FIXED} constraint
        (\texttt{attach\_ball()}) rigidly couples the ball to the end-effector; the
        constraint is removed at the release timestep (\texttt{release\_ball()}).

  \item \textbf{Velocity override.}  Immediately after constraint removal,
        \texttt{resetBaseVelocity} sets the ball's linear velocity to
        $\mathbf{v}_{\mathrm{cmd}}$.  Post-release drag is applied manually each
        simulation step via \texttt{applyExternalForce} using Eq.~35.
\end{enumerate}

The \texttt{vel\_limit\_multiplier} parameter only governs whether
\texttt{plan\_throw()} raises a feasibility warning; it does not alter the ball's
actual trajectory, which is governed entirely by $\mathbf{v}_{\mathrm{cmd}}$.

\needdata{architecture diagram (two-layer block diagram showing GP/policy/simulator on
top layer, NumPy sim and PyBullet on bottom layer, with $\mathbf{v}_{\mathrm{cmd}}$
as the interface arrow)}
